\section{Readings for each day}
\label{sec:advent-daily-readings}

Begin with today's date. December 8, 21 and 24 have directions of their
own below. On other days find the most recent Advent Sunday and turn to
that week. Read the Scripture in the middle column and pray the verse
in the last one. All the words for prayer are in \hyperref[sec:advent-home-prayers]{Praying at home}.
If you have more time, stay with the short reading from a Father given
beside the Scripture. On a day you return to prayer after an absence,
begin with the day you have reached.

The Sunday rows give the readings of the four Advent Masses. The Ember
and feast entries say which of their readings you will hear at Mass.
The other weekday passages are chosen for prayer at home. You can use
them on a saint's day even when your parish celebrates the saint's Mass.

\subsection{Find the week in your calendar}

Look across from the day of the week on which Christmas falls. The first
three Advent weeks each have seven days. The last column shows how much
of the fourth week you will have before the Vigil on December 24.

\begin{longtable}{@{}lp{0.24\linewidth}p{0.40\linewidth}@{}}
\textbf{Christmas} & \textbf{First Sunday} & \textbf{Fourth week before the Vigil} \\
\hline
\endhead
Sunday & November 27 & Sunday, December 18, through Friday, December 23 \\
Monday & December 3 & Go from Ember Saturday to Sunday's Vigil \\
Tuesday & December 2 & Sunday, December 23, alone \\
Wednesday & December 1 & Sunday, December 22, and Monday \\
Thursday & November 30 & Sunday, December 21, through Tuesday \\
Friday & November 29 & Sunday, December 20, through Wednesday \\
Saturday & November 28 & Sunday, December 19, through Thursday \\
\end{longtable}

\subsection{When a date and a weekday meet}

\begin{description}
\item[December 8.] Read Luke 1:26--28, the Gospel of the Immaculate
  Conception. Gabriel greets Mary before she has a word to offer him.
  Pray her Magnificat (Luke 1:46--55), dwelling on the great things
  God has done for her. In \hyperref[sec:advent-mary-word]{Mary receives the Word}, read Irenaeus
  on Eve and Mary's trust (\emph{Against Heresies} 3.22.4), then
  \hyperref[sec:advent-immaculate]{The Mother prepared by grace} on the feast. The feast takes
  precedence even on the Second Sunday of Advent. When they coincide,
  the priest adds the Sunday's Collect after the feast's, asking God to
  stir up our hearts. The next reading at home is Monday of the second week.
\item[December 17--23.] Keep the day's Scripture and its reading from a
  Father in the morning. At evening take the Scripture, O Antiphon
  and Magnificat under that date in \hyperref[sec:advent-evenings]{The Magnificat and the seven O Antiphons}. If
  December 17 is a Wednesday, for
  example, the morning Gospel is the Annunciation on Ember Wednesday.
  In the evening read Wisdom 8:1 and ask Wisdom to teach you prudence.
\item[December 21.] If today is Sunday, keep the Fourth Sunday's
  readings and prayer. On a weekday the universal calendar keeps
  St Thomas the Apostle and commemorates the Advent weekday at his Mass.
  Read the feast Gospel, John 20:24--29, and
  make Thomas's answer in verse 28 your own. In his \emph{Third
  Theological Oration} 29.19, Gregory of Nazianzus puts the Lord's
  burial beside the resurrection. Thomas knows the risen Lord by his
  wounds. If the feast falls on Ember Friday or Saturday, you may
  still read the Ember passage at home. A Saturday Mass at which
  orders are conferred keeps the Ember formulary. In each case the
  evening prayer for home is O Dayspring.
\item[December 24.] Turn to \hyperref[sec:advent-daily-vigil]{December 24, the Vigil} below. Even
  on Sunday the Vigil takes the place of Advent IV without a
  commemoration of the Sunday. There is never an ordinary fourth-week
  Saturday to make up.
\end{description}

\subsection{First week}

Listen at Mass for the psalm at the entrance. When the same man lifts
his soul to God again at the Offertory he still needs help. Learn his
words well enough to use them away from the book. Psalm 24 here is
Psalm 25 in many Bibles. The prayers below use the Douay--Rheims numbering.

\begin{longtable}{@{}lp{0.45\linewidth}p{0.31\linewidth}@{}}
\textbf{Day} & \textbf{Read} & \textbf{Pray} \\
\hline
\endhead
Sunday & Romans 13:11--14a and Luke 21:25--33. Who tells the frightened disciples they may lift their heads? \href{https://www.newadvent.org/fathers/310115.htm}{Cyril of Jerusalem, \emph{Catechetical Lecture} 15.1}, on the two comings. & Psalm 24:1--2. Name the fear you brought to Mass. \\
Monday & Isaiah 1:10--18. Listen for the widow whose cause has gone unheard. Cyril, Lecture 15.2, on what the refiner's fire discloses. & Psalm 24:4b--5. Ask to see whose cause you could take up. \\
Tuesday & Isaiah 2:1--5. Look at what neighbours can make from the metal of their weapons. Cyril, Lecture 15.3, on the renewal of creation. & Romans 15:5--6. Remember someone with whom it costs you something to speak peaceably. \\
Wednesday & Isaiah 11:1--10. Whose case does the promised king hear? Cyril, Lecture 15.4, on what the disciples ask Jesus. & Psalm 24:4b--5. Bring a judgment you made before you had heard the other person. \\
Thursday & Psalm 24:1--10. Where does the psalmist turn when he remembers the sins of his youth? Cyril, Lecture 15.26, on what you have given to Christ in a neighbour. & Psalm 50:12. Ask for the beginning of a change you can make today. \\
Friday & Luke 21:34--36. Continue where the Sunday Gospel ended. What can weigh a heart down? \href{https://www.newadvent.org/fathers/110108.htm}{Augustine, \emph{Confessions} 8.12.28}, under the fig tree. & Psalm 69:2. Say it again when that care comes back to mind. \\
Saturday & 1 Thessalonians 5:1--11. Listen for the work Paul gives believers to do for one another. Augustine, \emph{Confessions} 8.12.29. He reads to the end of Romans 13:14, beyond the words heard last Sunday. & Romans 15:13 for someone you mean to encourage. \\
\end{longtable}

\subsection{Second week}

John sends his disciples to Jesus with a question. As the days pass,
listen to the voices of the people for whom the answer is good news.
St Paul's two prayers in Romans 15 let you bring another person into
your prayer even on a day when you pray alone.

\begin{longtable}{@{}lp{0.45\linewidth}p{0.31\linewidth}@{}}
\textbf{Day} & \textbf{Read} & \textbf{Pray} \\
\hline
\endhead
Sunday & Romans 15:4--13 and Matthew 11:2--10. Which signs can the messengers carry back to John in prison? \href{https://archive.org/details/romanbreviaryref01cathuoft/page/229/mode/1up}{Gregory the Great, Homily 6 on the Gospels, §1}, on John's question. & Romans 15:13. Pray for someone still waiting for good news. \\
Monday & Isaiah 35:1--10. What are the people whose knees have grown weak told about their God? Gregory, Homily 6.1. Read what he says about the people who saw Christ's miracles. & Psalm 69:2. Say it with someone worn out by illness in mind. \\
Tuesday & Isaiah 40:1--11. Watch how the shepherd carries the lambs. Gregory, Homily 6.1, the closing words on Christ's Passion. & Romans 15:5--6. Remember a person who has had to be patient with you. \\
Wednesday & Isaiah 42:1--9. What does the servant do with a reed already bruised? \href{https://archive.org/details/romanbreviaryref01cathuoft/page/230/mode/1up}{Gregory, Homily 6.2}, on the reed Jesus speaks of. & Psalm 24:4b--5 before speaking to someone who has been hurt. \\
Thursday & Isaiah 61:1--4. Name someone in mourning as you read the promise to mourners. \href{https://www.newadvent.org/fathers/200136.htm}{Chrysostom, Homily 36 on Matthew, §2}, on John's disciples and their jealousy. & Romans 15:13 for that person, even if the two of you can only sit in silence. \\
Friday & Luke 4:16--22. Find the place Jesus reads to the people who knew him as a boy. Chrysostom, Matthew 36.2, on the sick whom Jesus heals. & Mary's Magnificat, especially Luke 1:53. Remember a household short of food. \\
Saturday & Luke 7:18--28. Wait to hear what Jesus says about John after the messengers leave. Chrysostom, Matthew 36.3, on the Lamb of God. & Romans 15:5--6. Remember the person who helped you find your way to Christ. \\
\end{longtable}

\subsection{Third week}

At Gaudete Mass Paul tells the Philippians to rejoice, though he is
writing from prison. As you come to pray, remember one good for which
you can give thanks while you are still waiting for help. Wednesday,
Friday and Saturday are the Advent Ember Days. Consult the December 21
entry above if the feast of St Thomas occurs. At the beginning of the
week pray \hyperref[sec:advent-gaudete-collect]{the Third Sunday's Collect}. Come back to it as often as you need to ask
for light.

\begin{longtable}{@{}lp{0.45\linewidth}p{0.31\linewidth}@{}}
\textbf{Day} & \textbf{Read} & \textbf{Pray} \\
\hline
\endhead
Sunday & Philippians 4:4--7 and John 1:19--28. Whom have the men questioning John failed to recognize? \href{https://archive.org/details/romanbreviaryref01cathuoft/page/241/mode/1up}{Gregory, Homily 7 on the Gospels, §1}, on John's humility. & The Collect of the Third Sunday. Remember someone who needs light for a decision. \\
Monday & Philippians 1:3--14. What good has come to other people through Paul's imprisonment? \href{https://www.newadvent.org/fathers/230214.htm}{Chrysostom, Homily 14 on Philippians}, on 4:4. & The Magnificat, Luke 1:46--55. Name one kindness you received during a hard day. \\
Tuesday & John 14:21--27. What does Jesus leave with the people from whom he is going away? Chrysostom, Philippians 14, on the peace that keeps guard (4:7). & Romans 15:5--6 for someone from whom you are about to part. \\
Ember Wed. & Luke 1:26--38, the Ember Gospel. Hear the news of Elizabeth before Mary gives her answer. Ambrose, \emph{Commentary on Luke} 2.14--15, in \hyperref[sec:advent-ember-wednesday]{Ember Wednesday at St Mary Major}. & Mary's answer, Luke 1:38. Bring a duty you have freely accepted. \\
Thursday & Isaiah 45:5--8. What are the heavens and the earth each asked to do? \href{https://www.newadvent.org/fathers/360312.htm}{Leo, Sermon 12.3}, on a poor harvest. & The petition for dew in Isaiah 45:8. Ask for fruit from some good work that seems to have failed. \\
Ember Fri. & Luke 1:39--47, the Ember Gospel. Elizabeth has heard Mary's greeting. Whom does she recognize? \href{https://archive.org/details/romanbreviaryref01cathuoft/page/252/mode/1up}{Ambrose, \emph{Commentary on Luke} 2.19}, on Mary's haste. & The Magnificat, Luke 1:46--55. Say Mary's reply to the end. \\
Ember Sat. & Daniel 3:46--51, from the longer Ember Mass. What do the three men do together inside the fire? Both Ember forms give Luke 3:1--6. Leo, Sermon 12.4, on the vigil at St Peter's. & Psalm 69:2. Remember a priest and a sick person to whom he brings Communion. \\
\end{longtable}

\subsection{Fourth week}

Use only the days that fall before December 24. At Mass John calls
people out to the Jordan. If you have only the Sunday this year, let
the word heard there have its answer today. What good could you still
do for someone who has been waiting on you?

\begin{longtable}{@{}lp{0.45\linewidth}p{0.31\linewidth}@{}}
\textbf{Day} & \textbf{Read} & \textbf{Pray} \\
\hline
\endhead
Sunday & 1 Corinthians 4:1--5 and Luke 3:1--6. Where is John when the word of God comes to him? \href{https://www.newadvent.org/fathers/220111.htm}{Chrysostom, Homily 11 on First Corinthians, §1}, on ranking our neighbours. & The petition for dew in Isaiah 45:8. Bring to God the days in which you wanted to begin again. \\
Monday & Luke 3:7--14. Which of John's instructions could you carry out today? Chrysostom, First Corinthians 11.5, on the motives behind a good deed. & Psalm 50:12 before returning something you owe or sharing what you have. \\
Tuesday & Micah 5:2--5. See where God looks for a shepherd for his people. Gregory, Homily 8 on the Gospels, §1, in \hyperref[sec:advent-bethlehem]{The house of bread}. & Mary's Magnificat, especially Luke 1:52. Remember someone whose work seldom receives thanks. \\
Wednesday & Luke 1:67--79. Hear the father who has only lately recovered his voice address his child. Ambrose, \emph{Commentary on Luke} 2.14--15, on Zachary's question, in \hyperref[sec:advent-ember-wednesday]{Ember Wednesday at St Mary Major}. & Romans 15:5--6 for someone dear to you who is far away. \\
Thursday & Hebrews 1:1--6. Hear how God has chosen to speak to us. \href{https://www.tertullian.org/fathers/cyril_on_luke_01_sermons_01_11.htm}{Cyril of Alexandria, Sermon 1 on Luke 2:7}, on where the eternal Son was laid. & Mary's answer, Luke 1:38. Leave a little silence to receive what you have heard. \\
Friday & John 6:47--52. Listen to Jesus tell those who come to him what bread he gives. Cyril, Sermon 1 on Luke 2:7, on what Christ gives at his table. & Thomas's confession in John 20:28 when you think of receiving the Lord at Christmas. \\
\end{longtable}

\subsection{December 24, the Vigil}
\label{sec:advent-daily-vigil}

In the daytime read the Vigil's Epistle, Romans 1:1--6, and its Gospel,
Matthew 1:18--21. Suppose you had heard only what Joseph hears. What
do you know about the child from the name you are to give him?
The Vigil Collect asks for confidence when the Redeemer comes again as
Judge. Pray Psalm 24:1--2, entrusting yourself to this child. In
St Cyril of Jerusalem's \emph{Catechetical Lecture} 15.1, read about the
Lord wrapped in swaddling clothes and the Lord clothed in light.

Christmas begins at First Vespers. At the crib or before the Mass in
the Night read Luke 2:1--14. Now you can join the angels' praise in
verse 14, given in \hyperref[sec:advent-home-prayers]{Praying at home}. For a longer watch take
St Gregory of Nazianzus, \href{https://www.newadvent.org/fathers/310229.htm}{\emph{Third Theological Oration} 29.19}.
He sets the cloths in which Mary wrapped her son beside the cloths
he will leave behind in the tomb.
